% Copyright (c) 2026 Eugene Joseph M. Ragasa.
% SPDX-License-Identifier: Apache-2.0
% Author: Eugene Joseph M. Ragasa
% Formal proof completed: 2026-10-28
\documentclass[11pt]{article}
\usepackage[T1]{fontenc}\usepackage{amsmath,amssymb}\usepackage[margin=1in]{geometry}\usepackage[hidelinks]{hyperref}
\title{Proof Dossier: Gauge--Truncation Noncommutation}\author{Eugene Joseph M. Ragasa}\date{Formal proof completed: 2026-10-28}
\begin{document}\raggedright\maketitle
\section{Identity and status}
ID \texttt{BG-FOURIER-THM-066}; status \texttt{proved}. Lean declaration:
\par\noindent{\tiny\path{ProjectKoios.BandGauge.Fourier.pointwiseSwapConjugation_zeroModeFourierTruncation_ne}}.
\section{Statement and proof}
For the constant field $H=\operatorname{diag}(1,0)$ on $\mathbb Z/2\mathbb Z$ and gauge values $U(0)=I$, $U(1)=\left(\begin{smallmatrix}0&1\\1&0\end{smallmatrix}\right)$,
\[\mathcal T_{\{0\}}(\mathcal G_UH)\ne\mathcal G_U(\mathcal T_{\{0\}}H).\]
The left side is constant after zero-mode truncation. The right side takes the distinct projector values at momenta zero and one. Equality would force those projectors to coincide, contradicting \path{BG-FOURIER-THM-065}.
\section{Boundary and audit}
This is one exact finite counterexample. It establishes neither a general convolution theorem nor continuum locality, smoothness, numerical evidence, or runtime correctness. Axioms are \texttt{propext}, \texttt{Classical.choice}, and \texttt{Quot.sound}; no admissions. Lean \texttt{4.35.0}, Mathlib {\scriptsize\path{56a005855ea17e0b0958f73ec801ae140b0a1679}}; AI-assisted, kernel-checked, no independent expert review recorded.
\end{document}
